%Lab template made by Joshua Milas for the EEEE labs

\def \labnum	{\#4}
\def \tonames   {Tyler Nicholson}
\def \disptitle	{MOSFET Differential Pair with Active Load}
\def \datestart	{04-03-2015}
\def \dateend	{04-17-2015}

%Lab template made by Joshua Milas for the EEEE labs
%Modified highly by Chris Culpepper



\documentclass[12pt]{article}	%12 point font

\usepackage{amsmath}	%Provides features and environments to help writing math equations, such as the align environment
\usepackage[pdftex]{graphicx}	%provides some extra features for including graphics. Optional
\usepackage[margin=0.8in]{geometry}	%Set the margins of the paper
\usepackage{float}	%Provides the [H] for hard setting things. It means "Put this exactly her
%\usepackage{circuitikz}	%Used for circuit drawing
\usepackage[export]{adjustbox}	%Used for alligning the RIT image to the right
\usepackage[outdir=./]{epstopdf}	%Used to include .eps vector images
\usepackage{indentfirst} %So the first line of a section is indented
\usepackage{titlesec} %To reduce the space after section titles
\usepackage{caption} %For caption junk
\usepackage{siunitx} % Formats the units and values
\usepackage{pgfplotstable} % Generates table from .csv
\usepackage{commath} %Used for math functions

\usepackage{datatool} %For the importation of CSVs to tables
\usepackage{colortbl} %So every other row of the table con be gray
\usepackage{fancyhdr} %For header and footer functionalithy
\usepackage{lastpage} %To get the last page num for the footer.
\usepackage{xcolor} %For different colors. 

\definecolor{light-gray}{gray}{0.85} %The gray color used for every other row in the table 

%To get numbers into Engineering notation, use \num{thhe number here} 
\sisetup{exponent-product = \times ,round-mode = figures,%
	  round-precision = 5, scientific-notation = engineering}

%These are the caption declarations for equations and tables. 
\DeclareCaptionType{equationCaption}[][List of equations]
\DeclareCaptionType{tableCaption}[][List of Tables]
\captionsetup[equationCaption]{name=Equation}%{labelformat=Equation}
\captionsetup[tableCaption]{name=Table}%{labelformat=Equation}

%For the pictures, i think
\epstopdfsetup{outdir=./}

%To move the section titles to match the format. 
\titlespacing\section{0pt}{5pt}{-2em}


\newcommand{\tab}{\hspace*{1.5em}}	%Tab command
\newcommand{\newsection}[1]{\setlength{ 	%Make a section
	\section*{#1}
	\leftskip}{0em}
	\vspace{5ex}
	%\setlength{\leftskip}{1em}
}
%This pretty prints an equation. 
%1st arg: The label to be used for reference 
%2nd arg: The equation. Use Math Mode. 
%3rd arg: The caption text. 
\newcommand{\eq}[3]{
	\begin{equationCaption}[H]
		\begin{equation}
			#2
			\label{#1}
		\end{equation}
		\caption{#3}
	\end{equationCaption}}
%Adds an image, captions and labels. 
%1: Optional arguments to \includegraphics. Use to set image params. 
%2: The filename
%3: The caption
%4: The label for referential purposes. 
\newcommand{\nimg}[4][]{
	\begin{figure}[H]
		\centering
		\includegraphics[#1]{#2}
		\caption{#3}
		\label{#4}
	\end{figure}
}
%Adds a table by CSV. Adds caption, label. 
%1 The label
%2 The filename
%3 The caption

%\newcommand{\ntab2}[3]{
	%\DTLloaddb[autokeys,noheader]{#1}{#2}
	%\arrayrulecolor{black}
	%\begin{table}[H]
		%\centering
%% Work out the column alignments.
		%\def\colalign{}%
		%\dtlforeachkey(\theKey,\theCol,\theType,\theHead)\in{#1}\do
			%{\edef\colalign{\colalign l}}%
%% Begin the tabular environment.
		%\edef\dobegintabular{\noexpand\begin{tabular}{\colalign}}%
			%\dobegintabular
			%%\toprule

%% Do the header row.
			%\gdef\doamp{\gdef\doamp{&}}%
					%%}%
				%%{\doamp\bfseries column\theCol}%
%% Iterate through the data.
			%\\%\midrule
			%\\[-5ex]
			%\DTLforeach{#1}{}{%
				%\\\gdef\doamp{\gdef\doamp{&}}%
				%\\[-3ex]\DTLifoddrow{\rowcolor{white}}{\rowcolor{light-gray}}%

				%\DTLforeachkeyinrow{\thisValue}{
					%\doamp 
					%\DTLifeq{\dtlrownum}{1}{
						%%aa\textbf{\thisValue}
						%\dtlkey
					%}{
						%\DTLifeq{\dtlcol}{1}{
							%%\dtlrownum
							%\thisValue
						%}{
							%\thisValue}
						%}}%
					%}
						%%\\\bottomrule
			%\end{tabular}
		%\caption{#3}
		%\label{#1}
	%\end{table}
%}
\newcommand{\ntab}[3]{
\DTLloaddb[autokeys]{#1}{#2}
	\arrayrulecolor{black}
	\begin{table}[H]
		\centering
% Work out the column alignments.
		\def\colalign{}%
		\dtlforeachkey(\theKey,\theCol,\theType,\theHead)\in{#1}\do
		{\edef\colalign{\colalign l}}%
% Begin the tabular environment.
		\edef\dobegintabular{\noexpand\begin{tabular}{\colalign}}%
			\dobegintabular
			%\toprule

% Do the header row.
			\gdef\doamp{\gdef\doamp{&}}%
			\renewcommand{\dtlbeforecols}{|}
			\renewcommand{\dtlaftercols}{|}
			\dtlforeachkey(\theKey,\theCol,\theType,\theHead)\in{#1}\do
			{\doamp\bfseries \theHead}%
% Iterate through the data.
			\\%\midrule
			\\[-5ex]
			\DTLforeach{#1}{}{%
				\\\gdef\doamp{\gdef\doamp{&}}%
				\renewcommand{\dtlbeforecols}{a}
				\renewcommand{\dtlaftercols}{a}
				\renewcommand{\dtlbetweencols}{a}
				\\[-3ex]\DTLifoddrow{\rowcolor{light-gray}}{\rowcolor{white}}%

				\DTLforeachkeyinrow{\thisValue}{
					\doamp 
					\DTLifeq{\dtlcol}{1}{
						\thisValue
					}{
						\num{\thisValue}}
					}}%
						%\\\bottomrule
					\end{tabular}
					\caption{#3}
					\label{#1}
				\end{table}
			}
\begin{document}
%\maketitle
\pagestyle{fancy} %Sets up the page to fancy for the header. 
\fancyhf{} % sets both header and footer to nothing
\renewcommand{\headrulewidth}{0pt} %Used to remove the spacing in a header.
% your new footer definitions here
\renewcommand{\headsep}{0pt} %Removes extra headerspacing
\lfoot{EEEE381 Tech Memo Lab \labnum} %The left part of the footer
\cfoot{\dateend} %The due date
\rfoot{\thepage\ of \pageref{LastPage}} %The right footer

 %Starts the RIT logo
\begin{flushright}
\begin{figure}[H]
\includegraphics[height=9ex,right,natwidth=883,natheight=167]{../title.png}
\end{figure}
\end{flushright}
\noindent
\\[-3em] %Removes extra space. Jankily. 
 %Starts the title
\huge
\textbf{EEEE 381 - Electronics I \\[1ex] Techical Memorandum}\\
\normalsize

 %This is the to/from garbage. Should probably make the name and section numbers variable
\noindent
\begin{tabular}{ll}
\textbf{From:} &Christopher Culpepper (Computer Engineering)\\
\textbf{Partner:} &N/A\\
\textbf{To:} &To: \tonames\ | Section L3\\
\textbf{Date:} &Performed: \datestart; Due: \dateend\\
\textbf{Subject:} &Lab \labnum: \disptitle
\end{tabular}

\noindent
\rule{\textwidth}{.1pt}
%\\[-2em]




\hyphenation{MOSFET} %Makes MOSFET not hyphenated


\newsection{Abstract}
The purpose of this exercise was to utilize the previously designed current source to implement an amplifier circuit. 
The modified Wilson current source was used as a bias current source for an active load differential amplifier with a single output. 
Once the circuit was designed and implemented correctly, it exhibited a gain that was close to the calculated value. 

\newsection{Theory}

The circuit used for this exercise is shown in figure \ref{fig:cir}. 
\nimg[width=6in]{fig:cir}{circuit.png}{Basic Amplifier Circuit}
By grounding the gates of the two NMOS transistors in the amplifier circuit, and assuming an $I_{ref}$ of 4mA, the Q-point values $V_{GSQ1}$ and $V_{GSQ2}$ can be found through the application of equation \ref{eq:vgsq}.
$V_{GSQ1}$ and $V_{GSQ2}$ were calculated to be 1.96V. 
\eq{eq:vgsq}{
	I_D = \frac{1}{2}k(V_{GS}-V_t)^2
	}{MOSFET Current Equation}


%2nd question
In order to allow bias current to flow through both sides of the amplifier, a common mode voltage, $V_{COM}$ needs to be applied. 
To ensure that all transistors remain in saturation, a minimum $V_{COM}$ value needs to be found. This value is calculated through equation \ref{eq:vcom}. This value was found to be <TODO>. 


	
%3
An important parameter of an amplifier is its gain. In the differential amplifier, there are two types of gain, differential and common mode. Differential gain refers to the gain when there is a differential voltage applied. Common mode gain is the gain of a common mode signal, or the same voltage applied to both inputs in addition to the differential signal.

Common mode signals are most likely noise that affects both inputs. Because the differential amplifier takes a difference between its two inputs, it is relatively good at filtering out common mode signals. 
To calculate the differential gain, $A_d$, equation \ref{eq:ad} is used. Calculation of the common mode gain, $A_{cm}$, uses equation \ref{eq:acm}. 
\eq{eq:ad}{
	A_d = (frac{v_o}{v_{diff}}) = g_m(r_{o2}||r_{o4}||R_L)}{
	Equation for the Calculation of the Differential Gain}
\eq{eq:acm}{
	A_{CM} = (\frac{v_o}{v_{cm}}) \approx - \frac{1}{2g_{m3}R_o}}{
	Equation for the Calculation of the Common Mode Gain}
	
The ratio of differential gain to common mode gain is called the common mode rejection ratio. This is usually measured in decibels, and is calculated through equation \ref{eq:cmrr}.
\eq{eq:cmrr}{
	CMRR (dB) = 20log\abs{\frac{A_d}{A_{CM}}}}{
	Equation for the Calculation of the Common Mode Rejection Ratio}

<TODO> RESULTS

The MOSFET parameters previously calculated can also be found through the use of PSpice simulation software. The simulated Q-point values were found to be <TODO>. 
Adding an AC signal to one of the gates of the amplifier, in addition to the common mode voltage, enables the study of the amplifier over many voltages. The circuit used to accomplish this is given in figure \ref{fig:precir5.png}.
\nimg[width=6in]{fig:precir5}{precir5.png}{AC Differential Input Circuit}
A plot of the data gathered is shown in figure \ref{fig:pre5}.
%\nimg{fig:pre5}{pre5.png}{Plot of the Output and Input Voltage for Figure \ref{fig:precir5}. 
Using the input and output voltages gathered from the simulation, the differential gain was simulated to be <TODO>. 
Similarly to how the differential gain was calculated, the common mode gain can be calculated using a simulation. This was done using the circuit in figure \ref{fig:precir6}. The plot of the data gathered is given in figure \ref{fig:pre6}. 

\nimg[width=6in]{fig:precir6}{precir6.png}{Common Mode Gain Input Circuit}
The common mode gain was calculated to be <TODO>, giving a common mode rejection ratio of <TODO>. 


\newsection{Results}
\newsection{Conclusion}
Current sources are important for amplifying circuits. The more stable the source, the bettor quality of amplification. This exercise reinforced theory learned in class and was valuable in that it taught the use and troubleshooting of important circuits which will be used later in lab. 
Even though the results for the output resistance were off for the Wilson source, the results were consistent, suggesting a problem with the measurements, or the chip. 
The results for the simple source were consistent and close to the expected values. 


\end{document}
